\documentclass[10pt,a4paper]{article}
\usepackage{../shared/proposal}
\usepackage[utf8]{inputenc}\usepackage[T1]{fontenc}
\usepackage{amsmath,amssymb}\usepackage{listings}\usepackage{xcolor}
\input{../shared/lstlang}

\begin{document}
\sloppy

\proposalhead{Delegated Sessions}{\texttt{openSession} already takes the session's owner as its
first argument. The gateway passes its own address.}{LP-320}

\section*{The problem}

The hosted gateway multiplexes arbitrarily many Web2 callers onto sessions funded by one consumer
wallet, behind careful engineering: idle-session claiming under \texttt{FOR UPDATE SKIP LOCKED}, a
cheapest-bid walk with per-provider cool-off, spend fuses that ship inert and arm per environment.
Better fuses do not reach what is left: one balance is a liveness dependency for every caller at
once; per-caller accounting lives in \texttt{credits\_ledger} and Redis, which the chain cannot
correct, while \texttt{routed\_sessions} carries no \texttt{user\_id}; and no caller has a claim on
the session serving them. The address that staked, that gets refunded, and that
\texttt{getUserSessions} returns is yours.

\section*{Mechanism}

The diamond carries a \texttt{Delegation} facet initialised against the \texttt{delegate.xyz} v2
registry. \texttt{openSession(address user\_, \dots)} calls \texttt{\_validateDelegatee} with
\texttt{DELEGATION\_RULES\_SESSION} at \texttt{SessionRouter.sol}:90, then pulls stake from
\texttt{user\_} at :104: caller and owner are separate variables, and
\texttt{SessionRouter.test.ts}:208 already tests the split. Nothing on
chain has to change for a gateway to stop holding its users' money. One argument in the Go node
forecloses it: \texttt{session\_router.go}:61 and :78 pass \texttt{opts.From}, fed by
\texttt{OpenSessionByBidId} reading \texttt{s.GetMyAddress(ctx)} (\texttt{service.go}:1260).

So the gateway holds a delegation instead of custody. The consumer writes
\texttt{delegateContract(gateway, diamond, DELEGATION\_RULES\_SESSION, true)} to the registry
itself; the gateway opens with \texttt{user\_} set to that consumer. Stake leaves the consumer's
balance and every exit is addressed back to them, \texttt{\_rewardUserAfterClose} to
\texttt{session.user} and \texttt{withdrawUserStakes} to \texttt{user\_}. That is the loss bound:
pooled, a compromised gateway key yields the pool; delegated, the worst it does is burn a stipend
claim and day-lock that consumer's principal for under a day, since the contract offers it no
other address.

Two caps already hold in deployed code and revoke without the gateway's help: the consumer's
ERC-20 allowance to the diamond, and \texttt{maxSessionDuration}. A boolean
\texttt{checkDelegateForContract} conveys nothing else, which should be said plainly: no per-open
size limit, no expiry, no model restriction. The registry holds the first already, since \texttt{delegateERC20(to,
MOR, rights, amount)} publishes an upper bound readable through \texttt{checkDelegateForERC20};
expiry and model scope want a small facet, which is the second move.

\textbf{The key problem, head on.} A delegation needs a consumer key, which account-native users
lack, and that cost does not disappear. But the gateway knows which callers
have one: \texttt{wallet\_links} stores addresses proved by SIWE, linked today to aggregate staking
allocations, and that cohort moves with no new onboarding. For the rest the pool stays, much
smaller, with a sponsored delegation from a smart account as the route off it.

\begin{stake}{Morpheus gets}
\item Delegated consumers stake their own MOR and draw their own stipend, so one wallet's
      stake-days stop rationing everyone's throughput.
\item Who consumed what, on chain in \texttt{userSessions[user\_]} rather than in a ledger the
      chain cannot correct, which is why \texttt{\_daylock\_from\_close\_receipt} had to be written.
\end{stake}

\begin{stake}{Hanzo and Lux get}
\item The shape \texttt{api.hanzo.ai} needs, which is why we will write it either way.
\end{stake}

\section*{First step, and who takes it}

Thread an optional session owner through \texttt{OpenSessionByBidId} into
\texttt{session\_router.go}, defaulting to \texttt{opts.From} so existing deployments behave
identically. That is the whole node change, its contract tests exist, and Hanzo opens it against
\texttt{dev}; the gateway then opens delegated for its SIWE-linked users first, behind the fuses it
ships. Neither needs a facet upgrade or a 5-of-9 signature. Reading the ERC-20 cap in
\texttt{\_validateDelegatee} is the one contract change, small enough for a queued diamond cut.

\noindent\textbf{What this does not solve.} Delegation fixes who owns a session, not what the
session bought. A delegated consumer is defrauded by a provider substituting its backend exactly as
a pooled one is, which is LP-313.

\noindent{\color{muted}\footnotesize\sffamily LP-320, full paper available. Market is the diamond at \texttt{0x6aBE1d282f72B474E54527D93b979A4f64d3030a} on Base.}

\end{document}
